\section{A sharp singular acquisition and memory law}\label{sec:singular-memory}
An acquired support may change when a positive but arbitrarily small preparation mass becomes visible. The mean budget then charges a long occupation on that component, even though a bounded terminal metric sees only its small probability. We give an exact finite-state law for this effect. It is a sharpness consequence of the general theory, not a replacement for the experiment-level theorem.

\subsection{The raw experiment and its actual occupations}
Fix $p\in(0,1)$ and write $q=1-p$. The physical states are $A$ (ready), $S$ (success) and $T$ (trapped). There is one sampling action and reports $f,s$. From $A$ the kernel sends the state to $S$ with report $s$ with probability $p$, and leaves it at $A$ with report $f$ with probability $q$. States $S,T$ are absorbing, with reports $s,f$, respectively. The preparation is
\begin{equation}\label{eq:rare-preparation}
 x_0^b=(1-b)e_A+be_T,\qquad 0\le b\le1.
\end{equation}
The raw instruments are nonnegative matrices on $\R^3$ with mass functional $\one^\top$. Positivity here means preservation of the positive cone, not a strictly positive lower bound on every transition. The next success report distinguishes $A$ from $T$, and a protected current audit distinguishes $S$ from them. These tests and the constant test span the affine prediction coordinates. Their raw positive update closes in three coordinates; a Fisher or covariance geometry is not an input.

We use two protected terminal tasks, scored only after the machine has committed to a terminal label. The first has the loss $\one_{\{X_\tau\ne S\}}$ and a singleton decoder body. The second audits a binary variable $Y$ that equals one at $S$ and is a fresh fair bit at $A,T$; the forecast lies in $[0,1]$ and the loss is $(Y-c)^2$. These are two separately declared tasks on the same sampling experiment. The audit costs one additional physical call and is not an observation available to the stopping rule.

For the failure task, stopping immediately on any success report never increases loss and can only decrease the mean call count. Merge all terminal labels. A machine using at most $M$ total states then has at most $k=M-1$ live states. Under a sequence consisting entirely of failure reports, its routing generates an absorption count $T_0\ge1$. If $b>0$, this is an \emph{actually prepared} component and $\E T_0<\infty$ follows from the mean budget. On the ready component the first success time $G_p$ is independent geometric with parameter $p$, independent of the routing coins. Hence, with
\[
 m=\E T_0,\qquad \varphi=\E q^{T_0},
\]
the actual failure probability and mean acquisition are exactly
\begin{equation}\label{eq:rare-flow}
 F=b+(1-b)\varphi,\qquad
 \E_b\tau=\frac{1-F}{p}+bm.
\end{equation}
The quantities $bm$ and $(1-F)/p$ are the actual trapped and ready live occupations. They are not normalized conditional substitutes for the preparation law.

\begin{theorem}[Exact memory frontier at every mean budget]\label{thm:rare-memory}
Let $b\in(0,1)$, $N\ge1$ and $M\ge2$, and put $k=M-1$. Let $g_{k,p}$ be the explicit discrete clock function in \eqref{eq:clock-frontier}. There is a unique $m_*\ge1$ satisfying
\begin{equation}\label{eq:rare-root}
 \frac{(1-b)(1-g_{k,p}(m_*))}{p}+bm_*=N.
\end{equation}
Define
\begin{equation}\label{eq:rare-profile}
 F_{p,b}(N,M)=b+(1-b)g_{k,p}(m_*).
\end{equation}
The optimal autonomous failure risk, over all homogeneous Borel randomized programs using at most $M$ total states, is exactly $F_{p,b}(N,M)$. It is attained by a clock with at most $k$ live states and one terminal state, with actual mean exactly $N$.

For the binary squared-prediction task on the same raw kernel,
\begin{equation}\label{eq:rare-prediction}
 \frac14 F_{p,b}(N,M)
 \le\calR_{\aut,\mathrm{pred}}^b(N,M)
 \le\frac12 F_{p,b}(N,M)-\frac14F_{p,b}(N,M)^2.
\end{equation}
The upper is implementable using the same clock and a single forecast. When $M=2$ it is the exact prediction risk. All statements hold also in the robust form with an unknown fixed preparation parameter in $[0,b]$, using one program designed for the known endpoint $b$.
\end{theorem}
\begin{proof}
The left side of \eqref{eq:rare-root}, regarded as a function of $m$, is continuous and strictly increasing from one to infinity, since $g_{k,p}$ is strictly decreasing and $b>0$. This proves existence and uniqueness.

Consider any competitor for the failure task. After the harmless modifications preceding the theorem, its all-failure clock has at most $k$ phases. Lemma~\ref{lem:discrete-clock} and \eqref{eq:rare-flow} imply
\[
 F\ge b+(1-b)g_{k,p}(m),\qquad
 m\le\frac{pN-1+F}{pb}.
\]
Since $g_{k,p}$ is decreasing,
\begin{equation}\label{eq:rare-fixed-inequality}
 F\ge b+(1-b)g_{k,p}\!\left(\frac{pN-1+F}{pb}\right).
\end{equation}
The argument is at least one for a feasible competitor. The right side is strictly decreasing in $F$, and equality occurs at \eqref{eq:rare-profile}. Thus \eqref{eq:rare-fixed-inequality} gives $F\ge F_{p,b}(N,M)$ for every competitor. The clock attaining $g_{k,p}(m_*)$ realizes equality in both equations of \eqref{eq:rare-flow}, proving the exact value and its state accounting.

For the prediction task, condition on the true terminal physical state as well as the terminal label. The fresh audit variance is $1/4$ on failure and zero on success. Every forecast, including a randomized one, therefore has risk at least $F/4$. Its failure probability cannot be smaller than the failure optimum already proved. For the upper use one terminal forecast $c=1-F_{p,b}(N,M)/2$. The audit has unconditional success probability $1-F_{p,b}(N,M)/2$, so its variance is the upper side of \eqref{eq:rare-prediction}. If $M=2$, the terminal forecast is necessarily a single constant (randomization cannot improve squared loss). Its best risk for failure probability $F$ is $F/2-F^2/4$, an increasing function on $[0,1]$. Stopping immediately on success leaves $F$ unchanged and reduces the mean. The failure optimum consequently also minimizes the prediction risk.

For the robust assertion, a fixed all-failure clock has $F_\theta=\theta+(1-\theta)\varphi$, increasing in $\theta$. Its mean is
$(1-\theta)\E\min(T_0,G_p)+\theta\E T_0$, also increasing in $\theta$. The endpoint construction is therefore feasible throughout $[0,b]$. Its chosen forecast is at least $1/2$, and its squared risk is increasing in the failure probability. The endpoint gives the lower bound on any common program. These facts prove the robust exact failure assertion, the prediction comparison, and the exact $M=2$ assertion without access to the true preparation parameter.
\end{proof}

At $b=0$ the failure optimum is $\max\{1-pN,0\}$. For $1\le N<1/p$ a one-phase randomized failure stop attains it; for $N\ge1/p$ waiting for the first success does so. Under this nominal preparation the all-failure path has no positive prepared mass, and its counterfactual clock may never halt. This does not contradict the preceding proof, which uses $b>0$ to charge its mean.

\begin{corollary}[Memory-dependent singular exponent]\label{cor:memory-exponent}
At the critical mean budget $N=1/p$, fix $k=M-1$. If
\begin{equation}\label{eq:phase-range}
 b\bigl(kp-1+q^k\bigr)\le q^k,
\end{equation}
then $m_*\ge k$ and the exact failure risk is the unique $F\in[b,1]$ solving
\begin{equation}\label{eq:memory-polynomial}
 F=b+(1-b)\left(\frac{qkb}{F+qkb}\right)^k.
\end{equation}
For each fixed $p\in(0,1)$ and $M\ge2$, as $b\downarrow0$,
\begin{equation}\label{eq:memory-exponent}
 F_{p,b}(1/p,M)\sim(qk)^{k/(k+1)}b^{k/(k+1)}.
\end{equation}
Consequently the prediction risk is of this same order, with constants between $1/4$ and $1/2$ times the failure scale. The exponent depends on paid finite-state memory although the physical carrier always has three coordinates. The asymptotic equivalence is for fixed $p,M$, not uniformly over a growing number of states.
\end{corollary}
\begin{proof}
The mean of the deterministic $k$-call failure clock is
$(1-b)(1-q^k)/p+bk$. Comparing it with $1/p$ yields exactly \eqref{eq:phase-range}. Thus the second branch of \eqref{eq:clock-frontier} applies. The mean identity in \eqref{eq:rare-flow} gives $m_*=F/(pb)$, and substitution proves \eqref{eq:memory-polynomial}.

The root tends to zero: for each fixed $f>0$ the right side of \eqref{eq:memory-polynomial} tends to zero. On the other hand $F/b\to\infty$, since a bounded subsequence of $F/b$ would leave its right side bounded away from zero. Rewrite the equation as
\[
 (F-b)(F+qkb)^k=(1-b)(qkb)^k.
\]
Dividing by $F^{k+1}$ and using $b/F\to0$ proves
\eqref{eq:memory-exponent}. The prediction comparison follows from
\eqref{eq:rare-prediction}.
\end{proof}

\begin{theorem}[Sharp preparation perturbation with two total states]\label{thm:rare-preparation}
For $M=2$ and $N=1/p$, the exact failure risk is
\begin{equation}\label{eq:two-state-preparation}
 F=\frac{pb+\sqrt{p^2b^2+4qb}}2,
 \qquad \rho_* =\frac{pb}{F}.
\end{equation}
An optimal program halts on success, and halts after each failure with probability $\rho_*$. Its exact prediction risk is $F/2-F^2/4$. The nominal value at $b=0$ is zero. The preparation base-norm distance is $\delta_0=2b$, while the instruments and task are unchanged. Thus the value difference is of order $\sqrt{c_N\delta_0}$ as $b\downarrow0$, for fixed $N>1$. For $N\ge2$ and $0<b\le1/4$, the comparison constants can be chosen numerically, independently of $N$.
\end{theorem}
\begin{proof}
For $k=1$, \eqref{eq:memory-polynomial} reduces to
$F^2-pbF-qb=0$. The positive root is \eqref{eq:two-state-preparation}. The trapped occupation is $b/\rho_*$ and, at mean equality, equals $F/p$, proving the routing formula. Theorem~\ref{thm:rare-memory} proves its global optimality and the prediction identity, including the robust interval assertion. The two preparations differ by $b(e_T-e_A)$, of base norm $2b$. Finally $\sqrt{qb}\le F\le pb+\sqrt{qb}$ and $F/4\le F/2-F^2/4\le F/2$. For $N\ge2$, $q\ge1/2$, giving the stated uniform comparisons.
\end{proof}

\subsection{Sharp one-call perturbation, including acquisition scaling}
Keep the preparation at $A$. Alter only its failure transition: from $A$, report $f$ and stay at $A$ with probability $q-h$, report $f$ and enter $T$ with probability $h$, and report $s$ and enter $S$ with probability $p$. Here $0<h<q$; the other transitions are unchanged. The joint raw-instrument base-norm distance from $h=0$ is exactly $\eta=2h$. In particular this is a primitive error, not a difference between normalized posterior parametrizations.

\begin{theorem}[Sharp one-call sensitivity at the active budget]\label{thm:rare-instrument}
At $N=1/p$ and $M=2$, put $a=(N-1)h$. The optimal failure risk over all legal randomized programs is
\begin{equation}\label{eq:two-state-instrument}
 F=\frac{\sqrt{a^2+4a}-a}{2},\qquad
 \rho_* =\frac{h(1-F)}{F+h(1-F)}.
\end{equation}
The program halts on success and halts after a failure with probability $\rho_*$. Its mean equals $N$, and the prediction risk is exactly $F/2-F^2/4$. The same program and forecast solve the robust problem with unknown fixed $h'\in[0,h]$.

For $0<a\le1/4$, $\sqrt a/2\le F\le\sqrt a$. Consequently the prediction-risk difference from $h=0$ is of order
$\sqrt{(N-1)\eta}$ with numerical constants, uniformly for $N\ge2$ in this small-error range. This matches the order in Theorem~\ref{thm:normalized-transfer}; no strictly larger exponent of the one-call perturbation can hold uniformly on its full class at this state budget.
\end{theorem}
\begin{proof}
With two total states there is one live state and one terminal state. A program is therefore specified by its stopping probabilities on $s,f$ and its terminal forecast. Halting immediately on success preserves the terminal physical success event and decreases the mean. Since squared risk optimized over the one forecast is increasing in the failure probability, this reduction is valid for both tasks, and includes all competitors.

Write $\rho>0$ for the failure stopping probability and
$D=p+q\rho+h(1-\rho)$. Direct summation of the two live occupations gives
\begin{align}
 z_A&=D^{-1},&z_T&=\frac{(1-\rho)h}{\rho D},\notag\\
 T(\rho)&=\frac{1+h(1-\rho)/\rho}{D},&
 F(\rho)&=\frac{q\rho+h(1-\rho)}{D}.
 \label{eq:instrument-occupations}
\end{align}
The terminal masses at $S,A,T$ are $p/D$, $\rho(q-h)/D$, $h/D$, respectively, and sum to one. These identities use the actual unconditioned transition probabilities. As $\rho$ increases from zero to one, the mean decreases continuously from infinity to one and $F(\rho)$ increases strictly; the latter follows from $F'(\rho)=p(q-h)/D^2>0$. Mean monotonicity also follows by coupling the failure-stop coins, and is strict here. Hence the optimum has $T(\rho)=1/p$.

Using $1-F=p/D$, this equality is
$F=h(1-\rho)(1-F)/\rho$. Substitute it into $D$ to obtain
$F^2=a(1-F)$, proving both formulas in \eqref{eq:two-state-instrument}. The prediction identity follows as in Theorem~\ref{thm:rare-memory}. For a fixed $\rho$, both $T(\rho)$ and $F(\rho)$ are increasing in $h$; explicitly their derivatives in $h$ are $p(1-\rho)^2/(\rho D^2)$ and $p(1-\rho)/D^2$. The endpoint program is therefore feasible and least favourable throughout the unknown parameter interval.

The raw change in the failure instrument moves mass $h$ from $A$ to $T$, giving induced $\ell^1$ norm $2h$; other input columns and reports are unchanged. The inequality $F^2=a(1-F)$ gives $F\le\sqrt a$, while $a\le1/4$ implies $F\le1/2$ and hence $F\ge\sqrt{a/2}\ge\sqrt a/2$. The prediction-risk comparison with $F$ proves the last claim.
\end{proof}

\subsection{Finite descriptions without budget enlargement}
The exact frontier concerns Borel programs, and its optimal probabilities may be irrational. This is separate from finite description. For rational $p,b,N$ with $b>0$, the clock in Theorem~\ref{thm:rare-memory} has certified finite approximants on the same states. Indeed $m_*\in[1,N/b]$, and the increasing function in \eqref{eq:rare-root} is explicitly computable with rational values at rational arguments on each branch. Bisection gives an interval enclosing $m_*$ to width at most $v/2$. Choose its lower endpoint $m$. If a phase boundary $k$ lies in the interval, it may first be split there. If $m\le k$, round the last continuation probability downward; if $m\ge k$, round the phase-advance probability $k/m$ upward. Either operation decreases the all-failure clock in stochastic order, so it decreases the actual mean. No comparison of an undecidable numerical boundary is required.

To make the additional rounding change in $m$ at most $v/2$, a dyadic mesh $Q^{-1}$ is sufficient if
\[
 Q^{-1}\le (v/2)\min\{1,kb^2/N^2\}.
\]
On the second branch $\rho=k/m\ge kb/N$ and
$k/\rho-k/(\rho+Q^{-1})\le N^2/(kb^2Q)$; the first branch is immediate. Since $g_{k,p}$ is $p$-Lipschitz, a total reduction of $m$ by $v$ raises failure risk by at most $pv$. Approximating the chosen constant forecast to accuracy $v$ adds at most $v^2$ to its optimal squared risk. The table uses at most $M$ persistent labels, a stored probability and readout, and a finite list of phase successors. Program bits and temporary coin-generation workspace depend on $\log Q$ and the input description, and are not included for free in $M$. Sampling still uses at most mean $N$ calls; preparation and the terminal audit are additional physical operations. These statements give sufficient description and execution bounds, not a lower frontier for all computational resources.
